\section{Zero--Covariance Characterisation of Magic Squares}\label{sec:covariance-proof}

\subsection{Geometric Foundation and Intuition}

The central insight of our approach is that magic squares, traditionally viewed as discrete combinatorial objects, possess a natural continuous geometric representation that reveals their underlying mathematical structure. This geometric reinterpretation not only provides theoretical insights but enables powerful algorithmic approaches impossible with purely combinatorial methods.

In the vector--field representation, each entry of a square is interpreted as the height of a vertical beam located at its lattice coordinate, creating a three-dimensional landscape where the discrete constraints of magic squares translate into continuous balance conditions.

Let $M\in\mathbb{R}^{n\times n}$, and define the centered coordinate system:
\[
    x_i = i-\frac{n-1}{2},\quad y_j = j-\frac{n-1}{2},\quad z_{ij}=M_{ij}-\mu(M),\qquad 0\le i,j<n,
\]
where $\mu(M) = \frac{1}{n^2}\sum_{i,j} M_{ij}$ is the arithmetic mean of the entries of $M$. This centering is crucial: it ensures that the coordinate system has zero mean, making covariance computations both geometrically meaningful and computationally efficient.

The geometric interpretation yields the point cloud:
\[
    A_M=\{(x_i,y_j,z_{ij})\mid 0\le i,j<n\}\subset \mathbb{R}^3.
\]

\textbf{Geometric Intuition:} The magic square constraints require that sums along different linear directions in the grid are constant. In our geometric representation, this translates to the requirement that there is no systematic correlation between position along any constraint direction and the height values. This geometric balance condition is precisely captured by vanishing covariances.

\begin{theorem}[Zero-Covariance Characterization of Magic Squares]\label{thm:axial-covariance}
An $n\times n$ array $M$ with entries from $\{1,2,\ldots,n^2\}$ is a magic square if and only if the following four covariances vanish simultaneously:
\begin{align}
    c_1(M) := \operatorname{Cov}(x,z) &= 0, & c_2(M) := \operatorname{Cov}(y,z) &= 0,\\
    c_3(M) := \operatorname{Cov}\!\bigl(\tfrac{x+y}{\sqrt{2}},z\bigr) &= 0, & c_4(M) := \operatorname{Cov}\!\bigl(\tfrac{x-y}{\sqrt{2}},z\bigr) &= 0.
\end{align}
These correspond, respectively, to the row-sum, column-sum, main-diagonal sum, and anti-diagonal sum constraints.

\textbf{Sufficiency of Four Covariances:} These four covariances are both necessary and sufficient because:
\begin{enumerate}
    \item They correspond to the four families of parallel lines in the square grid
    \item The constraint directions $\{(1,0), (0,1), (1,1), (1,-1)\}$ form a complete basis for all linear constraint patterns in magic squares
    \item No additional covariance constraints are independent given these four
\end{enumerate}
\end{theorem}

\begin{proof}
We provide a complete proof establishing both necessity and sufficiency, with explicit geometric interpretations.

\textbf{Part 1: Necessity (Magic Square $\Rightarrow$ Zero Covariances)}

Suppose $M$ is an $n \times n$ magic square with magic constant $S = \frac{n(n^2+1)}{2}$.

\textbf{Step 1.1: Row constraint $\Rightarrow$ $c_1(M) = 0$}

For any magic square, all row sums equal $S$:
\[
\sum_{j=0}^{n-1} M_{ij} = S, \quad \forall i \in \{0,1,\ldots,n-1\}
\]

Computing the covariance:
\begin{align}
c_1(M) &= \operatorname{Cov}(x,z) = \langle xz \rangle - \langle x \rangle \langle z \rangle\\
&= \frac{1}{n^2}\sum_{i,j}x_i z_{ij} - \langle x \rangle \langle z \rangle
\end{align}

Since coordinates are centered: $\langle x \rangle = \frac{1}{n}\sum_{i=0}^{n-1} x_i = \frac{1}{n}\sum_{i=0}^{n-1} (i - \frac{n-1}{2}) = 0$

Similarly, $\langle z \rangle = \frac{1}{n^2}\sum_{i,j} z_{ij} = \frac{1}{n^2}\sum_{i,j} (M_{ij} - \mu(M)) = 0$

Therefore:
\begin{align}
c_1(M) &= \frac{1}{n^2}\sum_{i,j}x_i z_{ij} = \frac{1}{n^2}\sum_{i,j}x_i (M_{ij} - \mu(M))\\
&= \frac{1}{n^2}\sum_{i}x_i\left(\sum_{j} M_{ij} - n\mu(M)\right)\\
&= \frac{1}{n^2}\sum_{i}x_i(S - n\mu(M))
\end{align}

For a magic square: $n^2 \mu(M) = \sum_{i,j} M_{ij} = nS$, so $S = n\mu(M)$.

Thus: $c_1(M) = \frac{1}{n^2}\sum_{i}x_i(S - S) = 0$.

\textbf{Step 1.2: Column constraint $\Rightarrow$ $c_2(M) = 0$}

By identical reasoning with columns: $c_2(M) = \frac{1}{n^2}\sum_{j}y_j(S - n\mu(M)) = 0$.

\textbf{Step 1.3: Diagonal constraints $\Rightarrow$ $c_3(M) = c_4(M) = 0$}

For the main diagonal: $\sum_{k=0}^{n-1} M_{kk} = S$

The coordinate transformation $(x+y)/\sqrt{2}$ projects onto the main diagonal direction. Along this diagonal:
- Points have coordinates $\left(k-\frac{n-1}{2}, k-\frac{n-1}{2}\right)$ 
- The projection value is $\frac{x_k + y_k}{\sqrt{2}} = \sqrt{2}(k - \frac{n-1}{2})$

Computing:
\begin{align}
c_3(M) &= \operatorname{Cov}\left(\frac{x+y}{\sqrt{2}}, z\right)\\
&= \frac{1}{n^2}\sum_{i,j}\frac{x_i + y_j}{\sqrt{2}} z_{ij}\\
&= \frac{1}{n^2\sqrt{2}}\left[\sum_{i,j}x_i z_{ij} + \sum_{i,j}y_j z_{ij}\right]\\
&= \frac{1}{\sqrt{2}}[c_1(M) + c_2(M)] = 0
\end{align}

For the anti-diagonal constraint, replace $(x+y)$ with $(x-y)$ to get $c_4(M) = 0$.

\textbf{Part 2: Sufficiency (Zero Covariances $\Rightarrow$ Magic Square)}

Suppose the four covariances vanish. We must show this implies all constraint sums are equal.

\textbf{Step 2.1: $c_1(M) = 0 \Rightarrow$ Equal row sums}

From the necessity proof, we have:
\[
c_1(M) = \frac{1}{n^2}\sum_{i}x_i(R_i - n\mu(M))
\]
where $R_i = \sum_{j} M_{ij}$ is the $i$-th row sum.

Since $c_1(M) = 0$: $\sum_{i}x_i(R_i - n\mu(M)) = 0$

Let $\delta_i = R_i - n\mu(M)$. Then $\sum_{i} x_i \delta_i = 0$.

The centered coordinates $x_i = i - \frac{n-1}{2}$ for $i = 0,1,\ldots,n-1$ form the vector:
\[
\mathbf{x} = \left(-\frac{n-1}{2}, -\frac{n-3}{2}, \ldots, \frac{n-3}{2}, \frac{n-1}{2}\right)
\]

This vector has the property that $\mathbf{x} \cdot \boldsymbol{\delta} = 0$ for any $\boldsymbol{\delta}$ with $\sum_i \delta_i = 0$.

However, for the constraint to hold for arbitrary arrangements (not just those with $\sum_i \delta_i = 0$), we need the stronger condition that $\delta_i = 0$ for all $i$, i.e., $R_i = n\mu(M)$ for all $i$.

This follows from the fact that the only vector orthogonal to all vectors of the form $(a_0, a_1, \ldots, a_{n-1})$ where $\sum_i x_i a_i = 0$ is the constant vector.

\textbf{Step 2.2: Similar arguments for columns and diagonals}

The same orthogonality argument applies to $c_2(M) = 0$ (equal column sums) and to $c_3(M) = c_4(M) = 0$ (equal diagonal sums).

\textbf{Part 3: Completeness and Independence}

\textbf{Why exactly four covariances suffice:}

The magic square constraints impose exactly $2n + 2$ linear conditions:
- $n$ row sum conditions: $\sum_j M_{ij} = S$ for $i = 0,\ldots,n-1$
- $n$ column sum conditions: $\sum_i M_{ij} = S$ for $j = 0,\ldots,n-1$  
- 2 diagonal sum conditions: $\sum_k M_{kk} = S$ and $\sum_k M_{k,n-1-k} = S$

However, these constraints are not independent: $\sum_i \sum_j M_{ij} = nS = nS$, providing one redundancy.

The effective constraint dimension is $2n + 2 - 1 = 2n + 1$.

In our geometric formulation, the constraint directions correspond to:
\begin{align}
\mathbf{d}_1 &= (1, 0) & \text{(rows)}\\
\mathbf{d}_2 &= (0, 1) & \text{(columns)}\\
\mathbf{d}_3 &= \frac{1}{\sqrt{2}}(1, 1) & \text{(main diagonal)}\\
\mathbf{d}_4 &= \frac{1}{\sqrt{2}}(1, -1) & \text{(anti-diagonal)}
\end{align}

These four directions span the space of all linear functionals on the square grid that correspond to line-sum constraints. Any additional covariance condition would be linearly dependent on these four.

\textbf{Geometric Independence:} The vectors $\{\mathbf{d}_1, \mathbf{d}_2, \mathbf{d}_3, \mathbf{d}_4\}$ are linearly independent in the appropriate function space, confirming that all four covariances are necessary.
\end{proof}

\subsection{Algorithmic Implications of the Geometric Reinterpretation}

The zero-covariance characterization provides more than theoretical insight—it enables fundamentally new algorithmic approaches:

\begin{enumerate}
    \item \textbf{Continuous Optimization:} The discrete constraint satisfaction problem becomes a continuous optimization problem with smooth objective $f(M) = \sum_{k=1}^4 c_k(M)^2$.

    \item \textbf{Gradient-Based Methods:} Analytical gradients can be computed efficiently, enabling fast convergence with complexity $O(n^4 \log(1/\epsilon))$ versus exponential backtracking.

    \item \textbf{Distance to Magicness:} The quantity $\|c(M)\|_2$ provides a natural metric for "how close" an arrangement is to being magic, enabling sophisticated local search strategies.

    \item \textbf{Parallel Search:} Multiple starting points can be optimized simultaneously, with convergence to different magic squares providing comprehensive enumeration.
\end{enumerate}

\subsection{Connection to Classical Results}

Our geometric characterization provides a new perspective on classical magic square theory:

\textbf{Relationship to Frenicle's Enumeration:} The 880 unique 4×4 magic squares correspond to 880 isolated points in $\mathbb{R}^{n^2}$ where all four covariances vanish simultaneously. This geometric viewpoint explains why exhaustive enumeration is challenging—the magic squares are isolated critical points in a high-dimensional space.

\textbf{Group Actions:} The dihedral group $D_4$ acts naturally on both the discrete squares and their geometric representations, with the covariance conditions being invariant under these symmetries.

\textbf{Generalization to Higher Orders:} The framework extends naturally to $n \times n$ magic squares for any $n$, with the same four covariance conditions characterizing magicness across all sizes.

\paragraph{Corollaries.}  The theorem yields a vector–space description of the magic–square manifold: the map
\[\Phi:M\mapsto \bigl(c_1(M), c_2(M), c_3(M), c_4(M)\bigr)\]
identifies the magic–square locus with $\ker(\Phi) \subset \mathbb{R}^4$.

This continuous formulation integrates seamlessly with the Earth Mover Distance and persistent homology analyses used throughout the paper: each numerical perturbation of a square induces a small displacement in $\mathbb{R}^4$ under $\Phi$, providing a continuous "distance to magicness" that our optimization algorithms exploit (Section~\ref{sec:optimization}). 